% formålet med denne fila er at den gir oss et mindre rotete utganspunkt for fremtidige rapporter, da den ikke inneholder noe tekst eller figurer som kan bli plagiat eller være i veien. bare overskriftsstrukturen står igjen.

% kopier den inn i f.eks. oblig3 sånn her:
% mkdir -p oblig3/
% cp Redusert_mal_fra_odin/redusertmal.tex oblig3/oblig3.tex

% preamble
\documentclass[11pt,a4paper,twocolumn]{article}
\usepackage[norsk]{babel}
\usepackage[a4paper,
            top=2cm,
            bottom=2cm,
            left=3cm,
            right=3cm,
            marginparwidth=1.75cm]{geometry}
\usepackage{amsmath}
\usepackage[norsk]{cleveref}
\usepackage[shortlabels]{enumitem}
\usepackage{graphicx}
\usepackage[T1]{fontenc}
\usepackage[hyphens]{url}
\usepackage{lipsum} 
\usepackage{siunitx}
\sisetup{
  output-complex-root = \ensuremath{\mathit{j}},
  complex-root-position = before-number,
  exponent-product = \cdot,
  output-decimal-marker = {,}
}
\usepackage{tikz}
\usepackage{circuitikz}
\usetikzlibrary{arrows.meta,positioning,calc}
\usepackage{pgfplots}
\pgfplotsset{compat=1.18}
\setlength{\parindent}{0pt}

% metadata
\title{Rapportskriving i ELI1300 - Redusert mal}
\author{Odin Heggvold Bekkelund\\For Gruppe ADK57}
\date{31.08.2026} 

% definer en url til senere?
\urldef{\myoverleafurl}\url{https://tinyurl.com/2yfbtdez}

% selve dokumentet starter
\begin{document}
\maketitle % Setter sammen \title og \author og \date

\begin{abstract}
Sammendrag
\end{abstract}

% Innholdsfortegnelse automatisk generert:
%\tableofcontents

\section{Innledning}
Innledning

\subsection{Utstyrsliste}
innhold

\section{Teoretisk del}
innledning

\subsection{Første ting}
innhold

\subsection{Andre ting}
innhold

\subsection{Tredje ting}
innhold

\section{Praktisk arbeid}
innledning

\subsection{Første ting}
innhold

\subsection{Andre ting}
innhold

\subsection{Tredje ting}
innhold

\section{Resultater}
innledning

\section{Diskusjon}
innhold

\section{Konklusjon}
innhold

referansebruk: noe står i læreboka \cite{laerebok}

\begin{thebibliography}{9}
\bibitem{laerebok}
Nilsson og Riedel \textit{Electric Circuits}, 11. Utgave, Pearson.
\bibitem{analogtech}
\url{https://www.analogtechnologies.com/document/white_paper/AWPLED-04-LED-Forward-Voltage.pdf}
\end{thebibliography}

\end{document}
